% Frame 1: Title
% ------------------------------------------------------------------

% ------------------------------------------------------------------
% Frame 2: Today's Goal
% ------------------------------------------------------------------
\begin{frame}{今日のゴール / Today's Goal}
  \begin{block}{目標}
    \texttt{sf app new} で独自プロジェクトを作り、ネイティブ API で全センサにアクセスする
  \end{block}

  \vspace{0.5em}

  \begin{itemize}
    \item \texttt{sf app new} で独自ファームウェアプロジェクトを作成
    \item テンプレートの構造（\api{app\_main}, \api{ControlTask}）を理解
    \item \api{StampFlyState} で全センサ・推定値に直接アクセス
    \item Teleplot でセンサデータをリアルタイム可視化
  \end{itemize}
\end{frame}

% ------------------------------------------------------------------
% Frame 3: Workflow Change
% ------------------------------------------------------------------
\begin{frame}[fragile]{開発フロー / Development Flow}

  \begin{alertblock}{注意: これまでのレッスンとやり方が変わります}
    {\small
    ワークショップ（\texttt{firmware/workshop/}）の外に、独立したプロジェクトを作ります。
    }
  \end{alertblock}

  \vspace{0.3em}

  {\small
  \begin{tabular}{lll}
    \hline
    & \textbf{L0\,--\,L10（これまで）} & \textbf{L11（このレッスン）} \\
    \hline
    開始 & \code{sf lesson switch N}
         & \code{sf app new my\_drone} \\
    編集先 & \texttt{firmware/workshop/main/}
           & \texttt{firmware/my\_drone/main/} \\
    ファイル & \texttt{user\_code.cpp}
             & \texttt{main.cpp} \\
    ビルド & \code{sf lesson build}
           & \code{sf build my\_drone} \\
    書き込み & \code{sf lesson flash}
             & \code{sf flash my\_drone -m} \\
    API & ws:: ラッパー & ネイティブ（StampFlyState） \\
    \hline
  \end{tabular}
  }
\end{frame}

% ------------------------------------------------------------------
% Frame 4: Workshop -> Native API comparison
% ------------------------------------------------------------------
\begin{frame}[fragile]{API 比較 / API Comparison}
  \begin{columns}[t]
    \begin{column}{0.46\textwidth}
      \begin{exampleblock}{ws:: ワークショップ}
        {\scriptsize
        \begin{itemize}
          \item \api{ws::gyro\_x()}, \api{ws::alt()}
          \item \api{ws::print(">tag:\%.2f", v)}
          \item \api{setup()} / \api{loop\_400Hz(dt)}
          \item 1 関数 = 1 センサ値
        \end{itemize}
        }
      \end{exampleblock}
    \end{column}
    \begin{column}{0.46\textwidth}
      \begin{block}{ネイティブ開発}
        {\scriptsize
        \begin{itemize}
          \item \api{state.getIMUData(a, g)}
          \item \api{printf(">tag:\%.2f\textbackslash n", v)}
          \item \api{app\_main()} / \api{ControlTask()}
          \item 構造体で複数値を一括取得
        \end{itemize}
        }
      \end{block}
    \end{column}
  \end{columns}

  \vspace{0.5em}
  {\footnotesize
    ネイティブ = vehicle ファームウェアと同じ構成。あらゆるアルゴリズムを実装可能。
  }
\end{frame}

% ------------------------------------------------------------------
% Frame 5: sf app new
% ------------------------------------------------------------------
\begin{frame}[fragile]{プロジェクト作成 / Create Project}
  \begin{lstlisting}[style=sfbash]
sf app new my_drone    # firmware/my_drone/ が生成
sf build my_drone      # ビルド
sf flash my_drone -m   # 書き込み + モニタ
  \end{lstlisting}

  \vspace{0.5em}

  \begin{block}{テンプレートの特徴}
    \begin{itemize}
      \item vehicle と同じセンサタスク・初期化コードを再利用
      \item \api{ControlTask()} のみを独自実装
      \item IMU, 気圧, ToF, 光学フロー, モーター等すべて利用可能
    \end{itemize}
  \end{block}
\end{frame}

% ------------------------------------------------------------------
% Frame 6: Project File Structure
% ------------------------------------------------------------------
\begin{frame}[fragile]{プロジェクト構成 / Project Structure}
  \begin{lstlisting}[style=sfbash, basicstyle=\ttfamily\small]
firmware/my_drone/
  CMakeLists.txt        # ビルド設定
  main/
    CMakeLists.txt      # コンパイル対象の定義
    main.cpp            # ← ここを編集
  \end{lstlisting}

  \vspace{0.3em}

  {\footnotesize
  \begin{tabular}{ll}
    \hline
    \textbf{ファイル} & \textbf{役割} \\
    \hline
    \texttt{CMakeLists.txt}      & vehicle/components と common を参照（全コンポーネント利用可能） \\
    \texttt{main/CMakeLists.txt} & vehicle のタスク（IMU, Baro, ToF 等）と init.cpp を再利用 \\
    \texttt{main/main.cpp}       & ユーザーが編集する唯一のファイル（ControlTask + コールバック） \\
    \hline
  \end{tabular}
  }
\end{frame}

% ------------------------------------------------------------------
% Frame 7: main.cpp Overview
% ------------------------------------------------------------------
\begin{frame}{テンプレート構造 / Template Structure}

  \begin{block}{\texttt{app\_main()} — ブートシーケンス}
    {\small NVS → センサ初期化 → ESKF 起動 → タスク起動（自動生成済み）}
  \end{block}

  \vspace{0.3em}

  \begin{alertblock}{\texttt{ControlTask()} — 400\,Hz ユーザーコード \textbf{← ここを編集}}
    {\small センサ読み取り・制御計算・モーター出力を記述するメインループ}
  \end{alertblock}

  \vspace{0.3em}

  \begin{exampleblock}{コールバック関数}
    {\small
    \begin{itemize}
      \item \api{onButtonEvent()} — ボタン押下で ARM / DISARM
      \item \api{handleControlInput()} — コントローラからの操縦入力を処理
    \end{itemize}
    }
  \end{exampleblock}
\end{frame}

% ------------------------------------------------------------------
% Frame 8: ControlTask Detail
% ------------------------------------------------------------------
\begin{frame}[fragile]{ControlTask / 制御ループ}
  \begin{lstlisting}[style=sfcpp, basicstyle=\ttfamily\scriptsize]
void ControlTask(void* pvParameters) {
    while (true) {
        xSemaphoreTake(g_control_semaphore, ...);
        // --- 1. センサ読み取り ---
        state.getIMUData(accel, gyro);    // 400 Hz
        state.getAttitudeEuler(r, p, y);  // ESKF 推定値
        state.getBaroData(alt, p);        // 50 Hz
        state.getToFData(bot, fnt);       // 30 Hz
        // --- 2. 制御計算（ユーザー実装） ---
        // PID, モーターミキシング等
        // --- 3. モーター出力 ---
        // g_motor.setThrust(0, thrust);
        // --- 4. Teleplot 出力（50 Hz） ---
        if (tick % 8 == 0) printf(...);
    }
}
  \end{lstlisting}

  {\footnotesize
    セマフォで IMU タイマー（2500\,$\mu$s）と同期。各センサは独立タスクで取得済み。
  }
\end{frame}

% ------------------------------------------------------------------
% Frame 9: StampFlyState API Table
% ------------------------------------------------------------------
\begin{frame}{StampFlyState API}
  {\small
  \begin{tabular}{lll}
    \hline
    \textbf{メソッド} & \textbf{ソース} & \textbf{取得データ} \\
    \hline
    \texttt{getIMUData(a, g)}           & BMI270          & 加速度 + ジャイロ \\
    \texttt{getBaroData(alt, p)}        & BMP280          & 高度 [m], 気圧 [Pa] \\
    \texttt{getMagData(mag)}            & BMM150          & 磁気 x,y,z [$\mu$T] \\
    \texttt{getToFData(b, f)}           & VL53L3CX        & 下方/前方距離 [m] \\
    \texttt{getFlowData(vx, vy)}        & PMW3901         & 速度 vx,vy [m/s] \\
    \texttt{getAttitudeEuler(r, p, y)}  & ESKF            & roll, pitch, yaw [rad] \\
    \texttt{getPowerData(v, i)}         & INA3221         & 電圧 [V], 電流 [A] \\
    \hline
  \end{tabular}
  }

  \vspace{0.5em}
  {\footnotesize
    \api{StampFlyState::getInstance()} でシングルトンを取得し、各メソッドを呼び出す。
  }
\end{frame}

% ------------------------------------------------------------------
% Frame 10: Hands-on
% ------------------------------------------------------------------
\begin{frame}[fragile]{実習: 全センサ Teleplot 可視化 / Hands-on}
  \begin{lstlisting}[style=sfcpp, basicstyle=\ttfamily\scriptsize]
auto& state = stampfly::StampFlyState::getInstance();
stampfly::Vec3 accel, gyro;
state.getIMUData(accel, gyro);
float alt, p;
state.getBaroData(alt, p);
float roll, pitch, yaw;
state.getAttitudeEuler(roll, pitch, yaw);

static uint32_t tick = 0; tick++;
if (tick % 8 == 0) {  // 50 Hz
    printf(">baro_alt:%.2f\n", alt);
    printf(">roll_deg:%.1f\n", roll * 57.3f);
    printf(">gyro_x:%.3f\n", gyro.x);
}
  \end{lstlisting}

  {\footnotesize
    \textbf{手順:} \texttt{sf app new my\_viz} → \texttt{firmware/my\_viz/main/main.cpp} の
    ControlTask を編集 → \texttt{sf build my\_viz} → \texttt{sf flash my\_viz -m} → Teleplot
  }
\end{frame}

% ------------------------------------------------------------------
% Frame 11: Checkpoint
% ------------------------------------------------------------------
\begin{frame}{チェックポイント / Checkpoint}
  \begin{alertblock}{確認事項}
    \begin{itemize}
      \item[$\square$] \api{sf app new} でプロジェクトを作成しビルドできた
      \item[$\square$] プロジェクト構成（3 ファイルの役割）を理解した
      \item[$\square$] \api{app\_main()} と \api{ControlTask()} の役割を理解した
      \item[$\square$] StampFlyState で全センサ値を直接取得できた
      \item[$\square$] Teleplot でリアルタイムグラフを確認した
    \end{itemize}
  \end{alertblock}

  \vspace{0.5em}

  \begin{block}{次のレッスン}
    Lesson 12: Python SDK プログラム飛行
  \end{block}
\end{frame}
